% Options for packages loaded elsewhere
\PassOptionsToPackage{unicode}{hyperref}
\PassOptionsToPackage{hyphens}{url}
\PassOptionsToPackage{dvipsnames,svgnames,x11names}{xcolor}
\documentclass[
  10pt,
]{article}
\usepackage{xcolor}
\usepackage[margin=0.83in]{geometry}
\usepackage{amsmath,amssymb}
\setcounter{secnumdepth}{-\maxdimen} % remove section numbering
\usepackage{iftex}
\ifPDFTeX
  \usepackage[T1]{fontenc}
  \usepackage[utf8]{inputenc}
  \usepackage{textcomp} % provide euro and other symbols
\else % if luatex or xetex
  \usepackage{unicode-math} % this also loads fontspec
  \defaultfontfeatures{Scale=MatchLowercase}
  \defaultfontfeatures[\rmfamily]{Ligatures=TeX,Scale=1}
\fi
\usepackage{lmodern}
\ifPDFTeX\else
  % xetex/luatex font selection
\fi
% Use upquote if available, for straight quotes in verbatim environments
\IfFileExists{upquote.sty}{\usepackage{upquote}}{}
\IfFileExists{microtype.sty}{% use microtype if available
  \usepackage[]{microtype}
  \UseMicrotypeSet[protrusion]{basicmath} % disable protrusion for tt fonts
}{}
\makeatletter
\@ifundefined{KOMAClassName}{% if non-KOMA class
  \IfFileExists{parskip.sty}{%
    \usepackage{parskip}
  }{% else
    \setlength{\parindent}{0pt}
    \setlength{\parskip}{6pt plus 2pt minus 1pt}}
}{% if KOMA class
  \KOMAoptions{parskip=half}}
\makeatother
\usepackage{listings}
\newcommand{\passthrough}[1]{#1}
\lstset{defaultdialect=[5.3]Lua}
\lstset{defaultdialect=[x86masm]Assembler}
\usepackage{longtable,booktabs,array}
\usepackage{calc} % for calculating minipage widths
% Correct order of tables after \paragraph or \subparagraph
\usepackage{etoolbox}
\makeatletter
\patchcmd\longtable{\par}{\if@noskipsec\mbox{}\fi\par}{}{}
\makeatother
% Allow footnotes in longtable head/foot
\IfFileExists{footnotehyper.sty}{\usepackage{footnotehyper}}{\usepackage{footnote}}
\makesavenoteenv{longtable}
\setlength{\emergencystretch}{3em} % prevent overfull lines
\providecommand{\tightlist}{%
  \setlength{\itemsep}{0pt}\setlength{\parskip}{0pt}}
\usepackage{bookmark}
\IfFileExists{xurl.sty}{\usepackage{xurl}}{} % add URL line breaks if available
\urlstyle{same}
\hypersetup{
  colorlinks=true,
  linkcolor={Maroon},
  filecolor={Maroon},
  citecolor={Blue},
  urlcolor={Blue},
  pdfcreator={LaTeX via pandoc}}

\author{}
\date{}

\begin{document}

{
\hypersetup{linkcolor=}
\setcounter{tocdepth}{3}
\tableofcontents
}
\section{Destructuring Field Guide}\label{destructuring-field-guide}

\textbf{Name:} Furkan Kaan Ünal\\
\textbf{Student number:} 2024502225\\
\textbf{Course:} MIS 233 --- Web Based Application Programming, Fall
2026\\
\textbf{Project basis:} the provided Assignment 4 starter

\subsection{Run and reproducibility}\label{run-and-reproducibility}

\begin{lstlisting}[language=bash]
npm install
npm run all > output.txt
npm run check
node js/depth.mjs
\end{lstlisting}

This guide uses a course registration and event-room scheduling dataset.
Each of the twelve forms is implemented in JavaScript and TypeScript.
The supplied TypeScript configuration enables
\passthrough{\lstinline!strict!} and
\passthrough{\lstinline!noUncheckedIndexedAccess!}.
\passthrough{\lstinline!@ts-expect-error!} is used only for the
intentional failure in T5; P7 invalid examples are shown as comments.

\textbf{Verification note (14 November 2026):} The complete assignment
was independently run on GitHub Actions using Node.js 22, the actual npm
dependencies, React 19 and TypeScript.
\passthrough{\lstinline!npm install --no-audit --no-fund!},
\passthrough{\lstinline!npm run all!},
\passthrough{\lstinline!npm run check!} and
\passthrough{\lstinline!node js/depth.mjs!} all passed. The full
captured output is committed in \passthrough{\lstinline!output.txt!},
including both JavaScript and TypeScript React rendering. Earlier local
npm installation attempts were blocked by DNS; this was resolved by
using a genuine CI runner.
\href{https://github.com/furkankaanunal/mis233-assignment4/actions/runs/43533531180}{Successful
verification run}.

\textbf{Bonuses claimed:} Bonus 1: LaTeX; Bonus 2: Depth; Bonus 3: Real
code; Bonus 4: Git repository
(\href{https://github.com/furkankaanunal/mis233-assignment4}{published
source}, verified tests).

\subsection{Field guide: the twelve
forms}\label{field-guide-the-twelve-forms}

An \emph{array pattern} consumes iterator values in order, while an
\emph{object pattern} requests properties by name. In both cases, the
right-hand expression is evaluated before individual bindings receive
their values. A \emph{rest pattern}
(\passthrough{\lstinline!...remaining!}) \textbf{collects} values,
whereas a \emph{spread expression} (\passthrough{\lstinline!...values!})
\textbf{expands} them.

\subsubsection{Form 1: Array return by
position}\label{form-1-array-return-by-position}

\textbf{Syntax:}
\passthrough{\lstinline!const [min, max] = examRange(scores)!}

\textbf{Example} (\passthrough{\lstinline!js/arrays.mjs!}; reduced to
the relevant lines):

\begin{lstlisting}
function examRange(scores) { return [Math.min(...scores), Math.max(...scores)] }
const [minimum, maximum] = examRange([62, 91, 74])
console.log('F1:', minimum, maximum)
\end{lstlisting}

\textbf{Output (captured from the JS run):}

\begin{lstlisting}
F1: 62 91
\end{lstlisting}

\textbf{Use it when:} A function returns two values with a fixed
positional meaning.

\textbf{Watch out:} Swapping the order changes the meaning; the return
should be a tuple in TypeScript.

\textbf{Plain assignments equivalent:} Read the corresponding array
positions in order. This is a convenient mental model, not a substitute
for iterator side effects or default-evaluation timing.

\subsubsection{Form 2: Skip and collect
rest}\label{form-2-skip-and-collect-rest}

\textbf{Syntax:}
\passthrough{\lstinline!const [, second, ...remaining] = agenda!}

\textbf{Example} (\passthrough{\lstinline!js/arrays.mjs!}; reduced to
the relevant lines):

\begin{lstlisting}
const agenda = ['registration', 'workshop', 'networking', 'closing']
const [, second, ...remaining] = agenda
console.log('F2:', second, remaining)
\end{lstlisting}

\textbf{Output (captured from the JS run):}

\begin{lstlisting}
F2: workshop [ 'networking', 'closing' ]
\end{lstlisting}

\textbf{Use it when:} The first array element is irrelevant but the rest
must be retained.

\textbf{Watch out:} The hole skips one iterator value; rest must come
last. Rest is collecting, not spread.

\textbf{Plain assignments equivalent:} Read the corresponding array
positions in order. This is a convenient mental model, not a substitute
for iterator side effects or default-evaluation timing.

\subsubsection{Form 3: Swap with
assignment}\label{form-3-swap-with-assignment}

\textbf{Syntax:} \passthrough{\lstinline![a, b] = [b, a]!}

\textbf{Example} (\passthrough{\lstinline!js/arrays.mjs!}; reduced to
the relevant lines):

\begin{lstlisting}
let morningRoom = 'B101', afternoonRoom = 'B202'
;[morningRoom, afternoonRoom] = [afternoonRoom, morningRoom]
console.log('F3:', morningRoom, afternoonRoom)
\end{lstlisting}

\textbf{Output (captured from the JS run):}

\begin{lstlisting}
F3: B202 B101
\end{lstlisting}

\textbf{Use it when:} Two existing variables must exchange values
without a temporary name.

\textbf{Watch out:} Use assignment, not \passthrough{\lstinline!const!};
both right-hand values are evaluated before either target is assigned.

\textbf{Plain assignments equivalent:} Read the corresponding array
positions in order. This is a convenient mental model, not a substitute
for iterator side effects or default-evaluation timing.

\subsubsection{Form 4: Array defaults}\label{form-4-array-defaults}

\textbf{Syntax:} \passthrough{\lstinline!const [x = fallback] = source!}

\textbf{Example} (\passthrough{\lstinline!js/arrays.mjs!}; reduced to
the relevant lines):

\begin{lstlisting}
const [firstSeat = 'unassigned', secondSeat = 'unassigned'] = [undefined, 'A12']
console.log('F4:', firstSeat, secondSeat)
\end{lstlisting}

\textbf{Output (captured from the JS run):}

\begin{lstlisting}
F4: unassigned A12
\end{lstlisting}

\textbf{Use it when:} Array data might omit a position or explicitly
contain \passthrough{\lstinline!undefined!}.

\textbf{Watch out:} A default does not replace
\passthrough{\lstinline!null!}, \passthrough{\lstinline!false!}, or
\passthrough{\lstinline!0!}.

\textbf{Plain assignments equivalent:} Read the corresponding array
positions in order. This is a convenient mental model, not a substitute
for iterator side effects or default-evaluation timing.

\subsubsection{Form 5: Object properties by
name}\label{form-5-object-properties-by-name}

\textbf{Syntax:}
\passthrough{\lstinline!const \{ attendee, track \} = registration!}

\textbf{Example} (\passthrough{\lstinline!js/objects.mjs!}; reduced to
the relevant lines):

\begin{lstlisting}
const registration = { attendee: 'Deniz', track: 'web', price: 250 }
const { attendee, track } = registration
console.log('F5:', attendee, track)
\end{lstlisting}

\textbf{Output (captured from the JS run):}

\begin{lstlisting}
F5: Deniz web
\end{lstlisting}

\textbf{Use it when:} A record is best understood by named fields
instead of positional indexes.

\textbf{Watch out:} An absent field becomes
\passthrough{\lstinline!undefined!}; destructuring
\passthrough{\lstinline!null!} or \passthrough{\lstinline!undefined!} as
the whole object throws.

\textbf{Plain assignments equivalent:} Look up the corresponding object
properties by name. This is a convenient mental model, not a substitute
for iterator side effects or default-evaluation timing.

\subsubsection{Form 6: Rename a
property}\label{form-6-rename-a-property}

\textbf{Syntax:}
\passthrough{\lstinline!const \{ code: courseCode \} = registration!}

\textbf{Example} (\passthrough{\lstinline!js/objects.mjs!}; reduced to
the relevant lines):

\begin{lstlisting}
const registration = { code: 'MIS233' }
const { code: courseCode } = registration
console.log('F6:', courseCode)
\end{lstlisting}

\textbf{Output (captured from the JS run):}

\begin{lstlisting}
F6: MIS233
\end{lstlisting}

\textbf{Use it when:} An object field needs a more descriptive local
variable name.

\textbf{Watch out:} \passthrough{\lstinline!code!} does not become a new
local variable: \passthrough{\lstinline!courseCode!} is the binding.

\textbf{Plain assignments equivalent:} Look up the corresponding object
properties by name. This is a convenient mental model, not a substitute
for iterator side effects or default-evaluation timing.

\subsubsection{Form 7: Object default}\label{form-7-object-default}

\textbf{Syntax:}
\passthrough{\lstinline!const \{ currency = "TRY" \} = registration!}

\textbf{Example} (\passthrough{\lstinline!js/objects.mjs!}; reduced to
the relevant lines):

\begin{lstlisting}
const registration = { code: 'MIS233' }
const { currency = 'TRY' } = registration
console.log('F7:', currency)
\end{lstlisting}

\textbf{Output (captured from the JS run):}

\begin{lstlisting}
F7: TRY
\end{lstlisting}

\textbf{Use it when:} Optional configuration fields need a reasonable
value.

\textbf{Watch out:} The default activates only for an
\passthrough{\lstinline!undefined!} property, never for
\passthrough{\lstinline!null!}.

\textbf{Plain assignments equivalent:} Look up the corresponding object
properties by name. This is a convenient mental model, not a substitute
for iterator side effects or default-evaluation timing.

\subsubsection{Form 8: Object rest}\label{form-8-object-rest}

\textbf{Syntax:}
\passthrough{\lstinline!const \{ code, ...details \} = registration!}

\textbf{Example} (\passthrough{\lstinline!js/objects.mjs!}; reduced to
the relevant lines):

\begin{lstlisting}
const registration = { code: 'MIS233', attendee: 'Deniz', track: 'web', seat: 'C4', price: 250 }
const { code, ...details } = registration
console.log('F8:', code, details)
\end{lstlisting}

\textbf{Output (captured from the JS run):}

\begin{lstlisting}
F8: MIS233 { attendee: 'Deniz', track: 'web', seat: 'C4', price: 250 }
\end{lstlisting}

\textbf{Use it when:} Some fields are handled separately and the
remaining enumerable own properties are forwarded.

\textbf{Watch out:} The rest object is shallow: nested objects still
point at the same references.

\textbf{Plain assignments equivalent:} Look up the corresponding object
properties by name. This is a convenient mental model, not a substitute
for iterator side effects or default-evaluation timing.

\subsubsection{Form 9: Nested
object/array}\label{form-9-nested-objectarray}

\textbf{Syntax:}
\passthrough{\lstinline!const \{ room: \{ floor \}, slots: [firstSlot] \} = timetable!}

\textbf{Example} (\passthrough{\lstinline!js/objects.mjs!}; reduced to
the relevant lines):

\begin{lstlisting}
const timetable = { room: { floor: 3 }, slots: ['09:00', '11:00'] }
const { room: { floor }, slots: [firstSlot] } = timetable
console.log('F9:', floor, firstSlot)
\end{lstlisting}

\textbf{Output (captured from the JS run):}

\begin{lstlisting}
F9: 3 09:00
\end{lstlisting}

\textbf{Use it when:} A response combines named object fields with
ordered array data.

\textbf{Watch out:} Missing \passthrough{\lstinline!room!} or
\passthrough{\lstinline!slots!} causes an exception unless you provide
suitable defaults; \passthrough{\lstinline!room!} is not bound as a
variable.

\textbf{Plain assignments equivalent:} Look up the corresponding object
properties by name. This is a convenient mental model, not a substitute
for iterator side effects or default-evaluation timing.

\subsubsection{Form 10: Computed
property}\label{form-10-computed-property}

\textbf{Syntax:}
\passthrough{\lstinline!const \{ [chosenField]: selectedValue \} = registration!}

\textbf{Example} (\passthrough{\lstinline!js/objects.mjs!}; reduced to
the relevant lines):

\begin{lstlisting}
const registration = { price: 250 }
const chosenField = 'price'
const { [chosenField]: selectedValue } = registration
console.log('F10:', selectedValue)
\end{lstlisting}

\textbf{Output (captured from the JS run):}

\begin{lstlisting}
F10: 250
\end{lstlisting}

\textbf{Use it when:} The property selected depends on a previously
computed key.

\textbf{Watch out:} Square brackets evaluate the key expression; the
resulting local name is \passthrough{\lstinline!selectedValue!}.

\textbf{Plain assignments equivalent:} Look up the corresponding object
properties by name. This is a convenient mental model, not a substitute
for iterator side effects or default-evaluation timing.

\subsubsection{Form 11: Parameter
patterns}\label{form-11-parameter-patterns}

\textbf{Syntax:}
\passthrough{\lstinline!function badge(\{ attendee = "Guest" \} = \{\})!}

\textbf{Example} (\passthrough{\lstinline!js/params.mjs!}; reduced to
the relevant lines):

\begin{lstlisting}
function badge({ attendee = 'Guest', track = 'general' } = {}) { return `${attendee} / ${track}` }
console.log('F11:', badge({ attendee: 'Deniz' }), '|', badge())
\end{lstlisting}

\textbf{Output (captured from the JS run):}

\begin{lstlisting}
F11: Deniz / general | Guest / general
\end{lstlisting}

\textbf{Use it when:} Callers supply an options object with optional
properties.

\textbf{Watch out:} The outer \passthrough{\lstinline!= \{\}!} handles
an omitted argument; inner defaults handle missing fields. Neither
repairs a literal null argument.

\textbf{Plain assignments equivalent:} Bind the parameters or loop entry
values one by one. This is a convenient mental model, not a substitute
for iterator side effects or default-evaluation timing.

\subsubsection{Form 12: Loop entry
patterns}\label{form-12-loop-entry-patterns}

\textbf{Syntax:}
\passthrough{\lstinline!for (const [key, value] of Object.entries(record))!}

\textbf{Example} (\passthrough{\lstinline!js/params.mjs!}; reduced to
the relevant lines):

\begin{lstlisting}
const inventory = { pens: 12, folders: 3 }
for (const [name, count] of Object.entries(inventory)) console.log('F12-object:', name, count)
\end{lstlisting}

\textbf{Output (captured from the JS run):}

\begin{lstlisting}
F12-object: pens 12
F12-object: folders 3
\end{lstlisting}

\textbf{Use it when:} A loop needs both the key and the associated value
of every entry.

\textbf{Watch out:} \passthrough{\lstinline!for...of!} iterates entry
arrays; \passthrough{\lstinline!for...in!} would iterate property names
instead. Map iteration yields tuples too.

\textbf{Plain assignments equivalent:} Bind the parameters or loop entry
values one by one. This is a convenient mental model, not a substitute
for iterator side effects or default-evaluation timing.

\subsection{Five pitfalls: observed errors and
explanations}\label{five-pitfalls-observed-errors-and-explanations}

\subsubsection{Pitfall 1 --- Missing right-hand
object}\label{pitfall-1-missing-right-hand-object}

\passthrough{\lstinline!const \{id\} = undefined!} and
\passthrough{\lstinline!const \{id\} = null!} throw before the property
can be looked up. The captured Node errors are:

\begin{lstlisting}
Pitfall1: undefined TypeError Cannot destructure property 'id' of 'input' as it is undefined.
Pitfall1: null TypeError Cannot destructure property 'id' of 'input' as it is null.
\end{lstlisting}

The fix is to check the object or use an outer parameter default when
omission is valid. An inner property default alone does not help.

\subsubsection{Pitfall 2 --- Undefined is not
null}\label{pitfall-2-undefined-is-not-null}

\passthrough{\lstinline!const \{a = 'fallback', b = 'fallback'\} = \{a: undefined, b: null\}!}
prints \passthrough{\lstinline!Pitfall2: \{"a":"fallback","b":null\}!}.
Only \passthrough{\lstinline!undefined!} triggers the binding default,
so a deliberate null stays null.

\subsubsection{Pitfall 3 --- Assignment versus
block}\label{pitfall-3-assignment-versus-block}

An existing binding can be reassigned with
\passthrough{\lstinline!(\{room\} = \{room:'new'\})!}; this prints
\passthrough{\lstinline!Pitfall3: new!}. Removing the parentheses leads
to the observed
\passthrough{\lstinline!Pitfall3-error: SyntaxError Unexpected token '='!}
in the tested Function constructor. The parser reads an initial
\passthrough{\lstinline!\{!} as a block, not an object assignment
pattern.

\subsubsection{Pitfall 4 --- No deep copy}\label{pitfall-4-no-deep-copy}

After \passthrough{\lstinline!const \{ room: shared \} = original!},
changing \passthrough{\lstinline!shared.floor!} also changes
\passthrough{\lstinline!original.room.floor!}; the run prints
\passthrough{\lstinline!Pitfall4: 5!}. Destructuring assigns the nested
reference. Even \passthrough{\lstinline!\{...rest\}!} makes only a
shallow copy.

\subsubsection{Pitfall 5 --- Name versus
position}\label{pitfall-5-name-versus-position}

An object \passthrough{\lstinline!\{first:'A', second:'B'\}!} can be
written as \passthrough{\lstinline!\{second, first\}!} without changing
values; an array \passthrough{\lstinline![a,b]!} depends on order. The
run prints \passthrough{\lstinline!Pitfall5: A B A B!}. Naming fields
reduces accidental positional confusion.

\subsection{TypeScript-specific lessons
T1--T5}\label{typescript-specific-lessons-t1t5}

\textbf{T1 --- Rename is not a type.} In
\passthrough{\lstinline!const \{ building: campusBuilding \} = location!},
the colon means \emph{rename}, not ``building has type campusBuilding''.
Put the annotation after the whole pattern:
\passthrough{\lstinline!const \{building\}: \{building: string\} = location!}.
\passthrough{\lstinline!ts/types.ts!} illustrates this with a typed
\passthrough{\lstinline!location!}.

\textbf{T2 --- Tuple and \passthrough{\lstinline!as const!}.}
\passthrough{\lstinline![number, string]!} guarantees that index 0 is a
number and index 1 is a string.
\passthrough{\lstinline![3, 'rooms'] as const!} gives a readonly tuple
with literal types \passthrough{\lstinline!3!} and
\passthrough{\lstinline!'rooms'!}, not a broadly mutable
\passthrough{\lstinline!(string | number)[]!}.

\textbf{T3 --- \passthrough{\lstinline!noUncheckedIndexedAccess!}.} If
\passthrough{\lstinline!attendees: string[]!} could be empty,
\passthrough{\lstinline!const [first] = attendees!} is
\passthrough{\lstinline!string | undefined!}. The explicit
\passthrough{\lstinline!first ?? 'no attendees'!} fallback makes empty
inputs safe. A fixed nonempty tuple has a known first element.

\textbf{T4 --- Optional field defaults.}
\passthrough{\lstinline!language?: string!} usually yields
\passthrough{\lstinline!string | undefined!}.
\passthrough{\lstinline!const \{language = 'tr'\} = options!} yields a
\passthrough{\lstinline!string!}. If
\passthrough{\lstinline!display?: string | null!},
\passthrough{\lstinline!display = 'standard'!} can still be
\passthrough{\lstinline!null!}: the default does not remove null.

\textbf{T5 --- A real type mismatch example.}
\passthrough{\lstinline!const \{count: countAsText\}: \{count: string\} = \{count: 8\}!}
causes
\passthrough{\lstinline!TS2322: Type 'number' is not assignable to type 'string'.!}
The marked \passthrough{\lstinline!@ts-expect-error!} line in
\passthrough{\lstinline!ts/types.ts!} intentionally keeps the error
visible to the compiler while allowing the project to check. The correct
solution is to use \passthrough{\lstinline!\{count: number\}!} or
\passthrough{\lstinline!String(count)!} after extracting it, not to lie
about the type. The runtime still stores the number.

\subsection{React props and children:
P1--P7}\label{react-props-and-children-p1p7}

\textbf{P1 --- One props object, two access styles.}
\passthrough{\lstinline!CardA(props)!} accesses
\passthrough{\lstinline!props.title!} and
\passthrough{\lstinline!props.children!};
\passthrough{\lstinline!CardB(\{title, children\})!} takes the same
fields using form 5 in a function-parameter position (form 11). Both
render
\passthrough{\lstinline!<section><h2>Schedule</h2>Friday</section>!}.
React does not pass \passthrough{\lstinline!title!} and
\passthrough{\lstinline!children!} as separate positional arguments.

\textbf{P2 --- Children is a prop.}
\passthrough{\lstinline!<CardB title="Equivalent">text</CardB>!} and
\passthrough{\lstinline!<CardB title="Equivalent" children="text" />!}
have equivalent children values and render
\passthrough{\lstinline!<section><h2>Equivalent</h2>text</section>!}. Do
not supply both syntaxes together; nesting is clearer when the content
is JSX.

\textbf{P3 --- Child value categories.} Strings and numbers render as
text; an element renders its markup; an array of elements renders them
in sequence. \passthrough{\lstinline!undefined!},
\passthrough{\lstinline!null!}, and \passthrough{\lstinline!false!}
produce no visible content. A function is \textbf{not} itself a
renderable ReactNode. The JS experiment deliberately passes a bare
function to a host element to observe React's handling, while
\passthrough{\lstinline!RenderCallback!} invokes a render-prop function
and renders its returned node. TS defines a callback separately as
\passthrough{\lstinline!(value: string) => ReactNode!}; it deliberately
rejects a function assigned to \passthrough{\lstinline!ReactNode!}.

\textbf{P4 --- Children versus rest.}
\passthrough{\lstinline!(\{title, children, ...rest\})!} keeps two
fields local and forwards remaining valid DOM props to
\passthrough{\lstinline!<section \{...rest\}>!}. Spreading
\passthrough{\lstinline!rest!} first and then setting
\passthrough{\lstinline!data-role="explicit"!} produces the explicit
value. Reversing that order lets the caller's
\passthrough{\lstinline!data-role="incoming"!} win. Use this pattern
carefully: forwarding arbitrary data to the DOM can leak unintended
attributes.

\textbf{P5 --- Named slots versus a single children flow.}
\passthrough{\lstinline!Slots!} takes \passthrough{\lstinline!header!}
and \passthrough{\lstinline!footer!} as named ReactNode props and
\passthrough{\lstinline!children!} as body. This is preferable when
layout positions need to be enforced independently.
\passthrough{\lstinline!Flow!} uses children only, allowing the caller
to control header/main/footer structure; this is simpler for highly
flexible compositions but less constrained.

\textbf{P6 --- Default value and null.}
\passthrough{\lstinline!function Status(\{label = 'Pending'\})!}
substitutes \passthrough{\lstinline!'Pending'!} when
\passthrough{\lstinline!label!} is missing or undefined.
\passthrough{\lstinline!null!} bypasses the destructuring default. The
sample uses \passthrough{\lstinline!label ?? '(empty)'!} inside JSX to
visibly distinguish the two cases:
\passthrough{\lstinline!<span>Pending</span>!} versus
\passthrough{\lstinline!<span>(empty)</span>!}.

\textbf{P7 --- Types and rejected props.}
\passthrough{\lstinline!children?: ReactNode!} permits an optional
renderable child (but not a bare function).
\passthrough{\lstinline!PropsWithChildren<\{title: string\}>!} adds
optional children. \passthrough{\lstinline!ComponentProps<'button'>!}
obtains built-in HTML button props including
\passthrough{\lstinline!type!}, \passthrough{\lstinline!disabled!} and
event handlers; it rejects \passthrough{\lstinline!href!}. The
implementation uses \passthrough{\lstinline!ButtonProps!} with
\passthrough{\lstinline!emphasis?!} to demonstrate combining a custom
prop with standard attributes. Deliberate P7 errors are documented as
commented-out examples, leaving T5 as the only
\passthrough{\lstinline!@ts-expect-error!} suppression.

\subsubsection{React static markup: real captured
output}\label{react-static-markup-real-captured-output}

These HTML lines were captured from a successful GitHub Actions run of
the actual JS/TS components; see \passthrough{\lstinline!output.txt!}.

\begin{lstlisting}[language=HTML]
P1-A: <section><h2>Schedule</h2>Friday</section>
P1-B: <section><h2>Schedule</h2>Friday</section>
P2-nested: <section><h2>Equivalent</h2>text</section>
P2-attribute: <section><h2>Equivalent</h2>text</section>
P3-string: <div>hello</div>
P3-number: <div>17</div>
P3-element: <div><strong>bold</strong></div>
P3-several: <div><em>one</em><em>two</em></div>
P3-undefined: <div></div>
P3-null: <div></div>
P3-false: <div></div>
P3-function-called: <div><strong>called</strong></div>
P3-function-uninvoked: <div></div>
P4-explicit-last: <section data-role="explicit" id="desk"><h3>Desk</h3>A</section>
P4-spread-last: <section data-role="incoming" id="desk"><h3>Desk</h3>A</section>
P5-slots: <article><header><b>Plan</b></header><main>Body</main><footer><small>End</small></footer></article>
P5-flow: <article><header><b>Plan</b></header><main>Body</main><footer><small>End</small></footer></article>
P6-undefined: <span>Pending</span>
P6-null: <span>(empty)</span>
P7-button: <button type="button" disabled="" class="primary">Save</button>
\end{lstlisting}

These are the React 19 server-rendered results from the verified
workflow, not guessed markup.

\subsection{JavaScript versus TypeScript (at least eight
distinctions)}\label{javascript-versus-typescript-at-least-eight-distinctions}

\begin{longtable}[]{@{}
  >{\raggedright\arraybackslash}p{(\linewidth - 6\tabcolsep) * \real{0.2500}}
  >{\raggedright\arraybackslash}p{(\linewidth - 6\tabcolsep) * \real{0.2500}}
  >{\raggedright\arraybackslash}p{(\linewidth - 6\tabcolsep) * \real{0.2500}}
  >{\raggedright\arraybackslash}p{(\linewidth - 6\tabcolsep) * \real{0.2500}}@{}}
\toprule\noalign{}
\begin{minipage}[b]{\linewidth}\raggedright
Form/topic
\end{minipage} & \begin{minipage}[b]{\linewidth}\raggedright
Same JavaScript behavior
\end{minipage} & \begin{minipage}[b]{\linewidth}\raggedright
Added TS type information
\end{minipage} & \begin{minipage}[b]{\linewidth}\raggedright
Mistake TS catches
\end{minipage} \\
\midrule\noalign{}
\endhead
\bottomrule\noalign{}
\endlastfoot
1 & Position-based \passthrough{\lstinline![min,max]!} & Return tuple
\passthrough{\lstinline![number,number]!} & Wrong tuple item type \\
2 & Skip and collect rest & \passthrough{\lstinline!string[]!} item can
be \passthrough{\lstinline!undefined!} & Calling
\passthrough{\lstinline!.toUpperCase()!} on possibly missing item \\
3 & Swap targets & Bindings must be assignable types & Swapping
\passthrough{\lstinline!string!} and \passthrough{\lstinline!number!}
variables \\
4 & Defaults only on undefined &
\passthrough{\lstinline!string \\| undefined!} source & Treating
\passthrough{\lstinline!null!} as \passthrough{\lstinline!string!} \\
5 & Extract properties by name & \passthrough{\lstinline!Registration!}
interface/type & Accessing a property absent from type \\
6 & Renaming with colon & Annotation goes after whole pattern &
Accidentally treating rename as a type \\
7 & Property fallback & Optional
\passthrough{\lstinline!currency?: string!} & Using optional result as
always defined \\
8 & Collect remaining properties & Inferred
\passthrough{\lstinline!Omit!}-like shape & Reading extracted
\passthrough{\lstinline!code!} from \passthrough{\lstinline!details!} \\
9 & Nested fields & Shape for \passthrough{\lstinline!room!} and tuple
\passthrough{\lstinline!slots!} & Reading absent nested field \\
10 & Computed key evaluated & Key constrained to
\passthrough{\lstinline!keyof Registration!} & Unknown key used on typed
object \\
11 & Defaulted options & \passthrough{\lstinline!BadgeOptions!}
parameter annotation & Invalid option value type \\
12 & Iterate entry pairs & \passthrough{\lstinline!Map<string,number>!}
& Treating a Map value as a string \\
P7 & Children still a prop & \passthrough{\lstinline!ReactNode!},
\passthrough{\lstinline!PropsWithChildren!},
\passthrough{\lstinline!ComponentProps!} & Passing
\passthrough{\lstinline!href!} to button or missing title \\
\end{longtable}

Some errors above are representative type mistakes rather than compiler
diagnostics captured from this environment.

\subsection{Bonus 1: LaTeX}\label{bonus-1-latex}

The directory \passthrough{\lstinline!latex/!} contains the full guide
as \passthrough{\lstinline!Destructuring\_Field\_Guide.tex!} and a
generated PDF. The PDF is typeset from the same README content, with
headings, a linked table of contents, real tables and syntax-preserving
code listings. Build command:

\begin{lstlisting}[language=bash]
cd latex
pdflatex -interaction=nonstopmode -halt-on-error Destructuring_Field_Guide.tex
pdflatex -interaction=nonstopmode -halt-on-error Destructuring_Field_Guide.tex
\end{lstlisting}

\subsection{Bonus 2: Depth}\label{bonus-2-depth}

The runnable extra experiments in \passthrough{\lstinline!js/depth.mjs!}
cover: strings (D1), Set iteration with duplicate removal (D2),
generator/iterator consumption (D3), lazy default evaluation (D4), using
an earlier binding inside a later default (D5), destructuring an error
in catch (D6), assignment to object properties (D7), Map entry iterators
(D8), and shallow object rest (D9). Run
\passthrough{\lstinline!node js/depth.mjs!}.

\begin{lstlisting}
D1-string: M I
D2-set: web [ 'finance' ]
D3-generator: A1 C3
D4-lazy-default: given computed 1
D5-earlier-binding: 3 6
D6-catch: registration closed
D7-target: { x: 5, y: 9 }
D8-map-entry: East 4
D9-shallow-rest: 7 9
\end{lstlisting}

\textbf{Why these happen:} an array pattern asks for an iterator, not
necessarily an Array. A Set iterates its unique insertion-order values,
while a generator produces items lazily. Skipped iterator positions can
still advance an iterator. Default expressions execute only for
\passthrough{\lstinline!undefined!}; consequently D4 calls
\passthrough{\lstinline!fallback()!} once rather than twice. A later
default can read an earlier initialized variable
(\passthrough{\lstinline!doubled = base * 2!}), but not a later variable
still inside its temporal dead zone. D6 destructures the thrown Error
object. D7 shows destructuring assignment can target properties. D9
proves an object rest copy is shallow.

\textbf{Readability tradeoff:}
\passthrough{\lstinline!const \{student:\{profile:\{address:\{city\}\}\}\} = record!}
is compact but fragile when any intermediate value is absent. Extracting
\passthrough{\lstinline!const city = record.student?.profile?.address?.city ?? 'unknown'!}
explicitly handles incomplete data and makes the fallback visible. For
short stable record shapes, destructuring remains clearer.

\textbf{Type-level extension:} \passthrough{\lstinline!ts/types.ts!}
defines \passthrough{\lstinline!First<T>!} and
\passthrough{\lstinline!Tail<T>!} with conditional types and
\passthrough{\lstinline!infer!} to extract tuple element types. This is
about types at compilation, not runtime destructuring.

\subsection{Bonus 3: Real code}\label{bonus-3-real-code}

Five source-grounded cases follow, using the \textbf{original supplied
starter} as real code written by someone else. Its line numbers are
stable because an untouched copy of the starter was provided with the
assignment. The external p-limit example is a sixth, additional case
(upstream line numbers are revision-dependent).

\begin{longtable}[]{@{}
  >{\raggedright\arraybackslash}p{(\linewidth - 8\tabcolsep) * \real{0.2000}}
  >{\raggedright\arraybackslash}p{(\linewidth - 8\tabcolsep) * \real{0.2000}}
  >{\raggedright\arraybackslash}p{(\linewidth - 8\tabcolsep) * \real{0.2000}}
  >{\raggedright\arraybackslash}p{(\linewidth - 8\tabcolsep) * \real{0.2000}}
  >{\raggedright\arraybackslash}p{(\linewidth - 8\tabcolsep) * \real{0.2000}}@{}}
\toprule\noalign{}
\begin{minipage}[b]{\linewidth}\raggedright
Case
\end{minipage} & \begin{minipage}[b]{\linewidth}\raggedright
Original source, exact line
\end{minipage} & \begin{minipage}[b]{\linewidth}\raggedright
Form
\end{minipage} & \begin{minipage}[b]{\linewidth}\raggedright
Equivalent without destructuring
\end{minipage} & \begin{minipage}[b]{\linewidth}\raggedright
Assessment
\end{minipage} \\
\midrule\noalign{}
\endhead
\bottomrule\noalign{}
\endlastfoot
1 & Starter \passthrough{\lstinline!js/arrays.mjs:8!} & 1
\passthrough{\lstinline![lowest, highest]!} &
\passthrough{\lstinline!const pair=minMax([4,9,2]); const lowest=pair[0]; const highest=pair[1]!}
& Clearly expresses a two-value return. \\
2 & Starter \passthrough{\lstinline!js/objects.mjs:6!} & 5
\passthrough{\lstinline!\{title,status\}!} &
\passthrough{\lstinline!const title=task.title; const status=task.status!}
& Short and readable because both properties are used. \\
3 & Starter \passthrough{\lstinline!js/params.mjs:5!} & 11
\passthrough{\lstinline!\{title,status='todo'\}=\{\}!} &
\passthrough{\lstinline!function createTask(options=\{\}) \{const title=options.title; const status=options.status===undefined?'todo':options.status; ...\}!}
& Makes the public configuration compact; nested defaults require
explanation. \\
4 & Starter \passthrough{\lstinline!js/props.jsx:10!} & 5 + 11
\passthrough{\lstinline!CardB(\{title,children\})!} &
\passthrough{\lstinline!function CardB(props)\{return <section><h2>\{props.title\}</h2>\{props.children\}</section>\}!}
& React-specific example: \passthrough{\lstinline!title!} is a named
attribute prop, \passthrough{\lstinline!children!} is the nested tag
content, but both are properties of the same object. \\
5 & Starter \passthrough{\lstinline!ts/props.tsx:9!} & 5 + 11
\passthrough{\lstinline!Card(\{title,children\}:CardProps)!} &
\passthrough{\lstinline!function Card(props: CardProps)\{return <section><h2>\{props.title\}</h2>\{props.children\}</section>\}!}
& Shows that TypeScript annotates the \textbf{entire parameter pattern},
not each renamed field. \\
6 &
\href{https://github.com/sindresorhus/p-limit/blob/main/index.js\#L7}{\passthrough{\lstinline!p-limit/index.js:7!}},
options branch & 5 + 7
\passthrough{\lstinline!(\{concurrency,rejectOnClear=false\}=concurrency)!}
&
\passthrough{\lstinline!const options=concurrency; concurrency=options.concurrency; rejectOnClear=options.rejectOnClear===undefined?false:options.rejectOnClear!}
& \textbf{Readability concern:} the expression reassigns
\passthrough{\lstinline!concurrency!} from an object to a number while
destructuring it. A separate \passthrough{\lstinline!options!} variable
would communicate intent more directly. The linked upstream source was
checked on 14 November 2026. \\
\end{longtable}

A strong improvement would be to add cases from the Assignment 3 Task
Manager and a locally installed library. Those sources were not supplied
in this upload, so their exact source lines have not been invented.

\subsection{Bonus 4: Git repository}\label{bonus-4-git-repository}

Public repository:
\textbf{https://github.com/furkankaanunal/mis233-assignment4}. The
project source, \passthrough{\lstinline!package-lock.json!}, captured
execution output, and strict TypeScript verification workflow are
published. See the
\href{https://github.com/furkankaanunal/mis233-assignment4/actions/runs/43533531180}{successful
GitHub Actions run}. The \passthrough{\lstinline!.gitignore!} excludes
\passthrough{\lstinline!node\_modules!}.

A local six-commit development history is packaged in
\passthrough{\lstinline!assignment4-history.bundle!} in the Moodle ZIP.
The GitHub repository also contains a visible incremental commit
history. Commits produced in one session were not falsely backdated to
imitate work over several days. Instructor access follows the course's
Moodle instructions.

\subsection{My own log}\label{my-own-log}

This section records verifiable tool-run outcomes, not invented personal
study experiences. Add your own wrong guesses or reflections before
submitting.

\begin{enumerate}
\def\labelenumi{\arabic{enumi}.}
\tightlist
\item
  The dependency install with
  \passthrough{\lstinline!npm ci --offline --no-audit --no-fund!}
  stopped with \passthrough{\lstinline!ENOTCACHED!} because the
  requested \passthrough{\lstinline!undici-types-6.21.0.tgz!} was not
  cached. \textbf{Learning:} lockfiles do not make an offline install
  possible without cached packages.
\item
  An online \passthrough{\lstinline!npm ci --no-audit --no-fund!}
  attempt timed out;
  \passthrough{\lstinline!curl -I https://registry.npmjs.org/react!}
  returned
  \passthrough{\lstinline!curl: (6) Could not resolve host: registry.npmjs.org!}.
  \textbf{Learning:} a partial \passthrough{\lstinline!node\_modules!}
  directory does not mean the dependencies are installed.
\item
  The JS null/undefined test produced two different TypeError
  descriptions. \textbf{Learning:} both fail as whole-object
  destructuring inputs; a property-level fallback cannot help.
\item
  The deliberately malformed assignment raised
  \passthrough{\lstinline!SyntaxError Unexpected token '='!}; wrapping
  the expression in parentheses succeeded. \textbf{Learning:} the
  grammar distinguishes statement blocks from destructuring assignments.
\item
  The shallow reference experiment produced
  \passthrough{\lstinline!Pitfall4: 5!} after changing
  \passthrough{\lstinline!shared.floor!}. \textbf{Learning:} extracting
  nested objects does not deep-clone them.
\item
  The lazy default experiment produced a call counter of
  \passthrough{\lstinline!1!} rather than \passthrough{\lstinline!2!}.
  \textbf{Learning:} the fallback function executes only for the element
  that was undefined.
\item
  \passthrough{\lstinline!D9-shallow-rest: 7 9!} showed the rest
  object's nested property change also affected the source.
  \textbf{Learning:} object rest is not a recursive copy.
\item
  \passthrough{\lstinline!tsc --noEmit!} initially reported
  \passthrough{\lstinline!TS17008: JSX element 'article' has no corresponding closing tag!}
  in \passthrough{\lstinline!ts/props.tsx!}; I added
  \passthrough{\lstinline!</article>!}, then syntax parsing succeeded.
  \textbf{Learning:} JSX tags must be balanced before meaningful
  type-checking can proceed.
\item
  \passthrough{\lstinline!tsc -p tsconfig.offline.json!} passed with no
  diagnostics for non-React TypeScript files. \textbf{Learning:} a
  targeted compile is helpful but does not replace a complete project
  check.
\item
  On GitHub Actions, the real \passthrough{\lstinline!npm install!},
  \passthrough{\lstinline!npm run all!},
  \passthrough{\lstinline!npm run check!}, and depth experiment
  completed successfully. \textbf{Learning:} a reproducible clean runner
  can validate the complete project even when the original machine has
  DNS restrictions.
\end{enumerate}

\subsubsection{References}\label{references}

\begin{itemize}
\tightlist
\item
  ECMA-262, destructuring binding and assignment patterns:
  https://tc39.es/ecma262/
\item
  MDN, destructuring assignment:
  https://developer.mozilla.org/en-US/docs/Web/JavaScript/Reference/Operators/Destructuring
\item
  TypeScript handbook, object types and tuples:
  https://www.typescriptlang.org/docs/handbook/2/objects.html
\item
  React documentation, passing props:
  https://react.dev/learn/passing-props-to-a-component
\item
  Third-party example source: https://github.com/sindresorhus/p-limit
\end{itemize}

\end{document}
